\section{System Design and Implementation}
\subsection{AI Model Integration}

\subsubsection{System Overview}
The system has three cooperating components: the variational circuit backend, the telemetry engine, and the reinforcement-learning scheduler. During a training run the backend performs one epoch of parameter optimization and exposes telemetry (gradient variance, loss, depth, and related features). The telemetry engine assembles the sixteen-dimensional state and passes it to the scheduler. The scheduler returns one of the five actions; a mutation engine applies the action to the circuit -- adding, soft-masking, or freezing layers -- under a commit/rollback guard. The loop repeats for the configured number of epochs. Algorithm~1 (Section~4.7.4) gives this loop formally; the description here focuses on which software component owns which responsibility, rather than on the control-flow logic itself.

\subsubsection{Component Responsibilities}
Each of the three components has a single, narrow responsibility, which is what allows the same scheduler and telemetry code to run unmodified against either physics backend (Section~4.7.1).

\begin{itemize}
\item Circuit backend (simulate or PennyLane): owns the quantum state, the parameter-optimization step, and the raw measurement outputs. It has no knowledge of the RL agent, the reward, or the action space; its only external contract is to accept one epoch's worth of optimization steps and return the resulting loss, accuracy, and gradient statistics.
\item Telemetry engine: owns the sixteen-dimensional state vector. It reads the raw statistics from the circuit backend, computes derived quantities ($\Delta$loss, the near-zero-gradient ratio, rel\_weak), and normalizes everything into the fixed-size vector the scheduler expects, regardless of which backend produced the raw numbers.
\item Reinforcement-learning scheduler and mutation engine: own the policy and the architecture-change logic. The scheduler consumes only the telemetry vector and returns only one of the five discrete actions (\Cref{tab:actions}); the mutation engine is the sole component permitted to modify the circuit's structure, and it does so only through the commit/rollback protocol of Algorithm~1.
\end{itemize}

\subsection{Data Flow and Processing}

\subsubsection{Preprocessing and Training Flow}
Raw data enters as CSV. The preprocessing pipeline first splits the data deterministically (by seed) into train and validation (70/30), then fits a standardization (Z-score) on the training split only, reconciles the feature count with the qubit count by PCA, one-to-one mapping, or zero-padding as appropriate (Section~4.2), and finally rescales features to the angle-encoding range [0, π]. The training split drives parameter optimization and the validation split drives the scheduler's accept/rollback decisions; the deployment and comparison scripts behind Section~6 do not additionally carve out a held-out test split (Section~4.5), though the shared loader supports one and it is used separately by the prediction-export script behind the demonstration interface (Section~5.3).

\subsubsection{Telemetry and Logging Flow}
Telemetry produced each epoch flows to the scheduler, and the resulting per-epoch record -- loss, accuracy, depth, action, commit/rollback counters -- is written to a deployment log that is the single source for every reported number.

\subsection{Deployment Strategy}

\subsubsection{Train-on-Surrogate, Deploy-on-PennyLane Workflow}
The confirmed workflow is train-on-surrogate, deploy-on-PennyLane. The RL policy is trained against the surrogate and then executed on the statevector simulator, which produces the deployment logs and the final trained circuit parameters used for inference. For the demonstration, the trained circuit is applied to held-out test samples to produce live single-sample predictions; because the parameters are captured immediately after training, no separate weight-persistence step is required for the demonstration to be reproducible.

\subsubsection{Reuse Across Experimental Regimes}
The same deployment procedure is reused, unmodified, for every regime reported in Section~6: only the four environment parameters described in Section~4.7.1 (dataset, qubit count, cost type, depth bounds) and the choice of physics backend change between the PennyLane-real run (\Cref{tab:heart}) and the simulate-backend run (\Cref{tab:sim8g}). This is a direct consequence of the component-responsibility split in Section~5.1 -- because the scheduler and telemetry engine never reference the backend directly, deploying against a new regime is a configuration change, not a code change, which is also why the regimes in Section~6 could be produced from two scripts (Appendix~C) rather than separate ones.

\subsection{Scalability and Maintenance}

\subsubsection{Scalability Axes}
The design scales along three axes. New datasets are added through a small preparation script and a single loader branch, without touching the model. New qubit counts or cost types are supported by re-calibrating the surrogate. Maintenance is simplified by the single-source-of-truth ansatz: any change to the circuit is made in one place and is automatically reflected in measurement, training, and deployment. The two-simulator split also means that expensive statevector evaluation is only incurred at deployment, keeping the iteration loop fast.

\subsubsection{Reuse in the Section 6 Sweep}
This same scalability axis is what let the Section~6.2 fixed-depth sweep be run as a batch of independent, differently-configured deployments rather than as bespoke code: the cost-type sweep is, mechanically, the same re-calibration-then-deploy procedure used for any other configuration change, repeated six times per regime for k = 2, 4, 6, 8, 12, 16. The maintenance cost of adding a new regime to Section~6 is therefore one surrogate re-calibration and one deployment run, not a new code path.

